\section{Conclusion}\color{black}
This paper presents a controlled short-window AMC study that separates
feature-geometry changes from network-capacity changes on a shared primary
backbone. Under the locked protocol, increasing width from $0.50\times$ to
$1.00\times$ improved final-epoch accuracy for both F2 and F7, while the
$1.00\times$--$2.00\times$ increase produced no distinguishable change for F2
and a decrease for F7. The feature contrast was not distinguishable at the two
lower widths and favored F2 at $2.00\times$; the interaction contrast was
$-0.68$ percentage points with Holm-adjusted $p=0.0312$.

These findings support a conditional conclusion about the tested feature ladder,
width range, 128-sample observation, two splits, and BB1 protocol. They do not
establish universal feature-geometry superiority, hardware latency or energy
advantages, or cross-corpus generalization. Core-FLOP-matched accuracy, SNR
region effects, and the secondary-backbone control remain outside the reported
results and are not inferred here.
